\documentclass{situationist-preprint}

\title{Necessary Truths: An Inventory}
\msid{TSM-11-b}
\msfamily{TSM-11}
\msdiscipline{Time, Space and Modality}
\mstype{Consequence}
\mskeywords{necessity; contingency; essence; laws of nature; identity; origin}

\begin{document}
\maketitle

\begin{abstract}
TSM-11-a argues that the necessary truths are exactly those of logic, mathematics and essence, with their logical consequences. We draw up the inventory this yields. Necessary: logical and mathematical truths, identities between things, truths of kind and constitutive origin, and the essential truths about properties and relations. Contingent: the existence of any concrete thing, the laws of nature, the values of the constants, the arrangement and number of things, and the existence of anything at all concrete. Each entry is derived from Theorem 2.3 of TSM-11-a and the account of essence of MET-13.
\end{abstract}

\section{The inventory}

By Theorem~2.3(i) of TSM-11-a, a truth is necessary exactly when it follows logically from the essential truths $E$; logic and mathematics are included because they follow from any premises or from the real definitions of mathematical objects. By Assumption~2.1 of TSM-11-a, real definitions state kind and constitutive origin, not location, accidental qualities, unrelated relations or laws. The entries below follow.

\paragraph{Necessary.}
\begin{enumerate}
\item \emph{Logic and mathematics}, including set theory so far as it follows from the definitions of sets.
\item \emph{Identities and distinctnesses.} If $a = b$, it is part of what $a$ is to be $a$, so $\nec(a = b)$; distinctness is necessary likewise \citep{Kripke1980}.
\item \emph{Kind.} A thing that is essentially a horse is necessarily a horse, if it exists.
\item \emph{Constitutive origin.} A thing originating from a given source originates from it in every world in which it exists \citep{Kripke1980}, where the origin belongs to its real definition (MET-13).
\item \emph{Essences of properties and relations}, such as that water is composed of H$_2$O molecules or that being red excludes being green \citep{Fine1994}.
\end{enumerate}

\paragraph{Contingent.}
\begin{enumerate}
\item \emph{Existence of concrete things.} Real definitions say what a thing is, not that it is; so \enquote{$a$ exists} does not follow from $E$, and by Theorem~2.3(ii) it is contingent.
\item \emph{Laws and constants.} No real definition includes the laws (Assumption~2.1 of TSM-11-a); the laws are therefore contingent, as is required by the best-system account of MET-12, and so are the values of the constants.
\item \emph{Location, quality and arrangement}, the number of things of a kind, and their spatiotemporal relations.
\item \emph{That there is something concrete.} The conjunction of the negations of all existence claims is consistent with $E$, since $E$ contains only conditional truths about what things are.
\end{enumerate}

\begin{objection}
If kinds are essential and laws flow from the natures of kinds, as dispositional essentialists hold \citep{Bird2007}, the laws are necessary.
\end{objection}

\begin{reply}
The inventory follows Assumption~2.1 of TSM-11-a, which excludes laws from real definitions. What the essence of a kind fixes is its composition and origin; how things of that kind behave in given circumstances is, on the account adopted here, a further fact. The dispositional essentialist must deny Assumption~2.1, and the inventory records what follows if it is accepted.
\end{reply}

\section{Conclusion}

The necessary truths form a short list: logic, mathematics, identity, kind, constitutive origin and the essences of properties. Everything that depends on which things exist, where they are, how they are arranged and what laws they obey is contingent.

\begin{thebibliography}{99}\small
\bibitem[Bird(2007)]{Bird2007} Bird, A. (2007). \emph{Nature's Metaphysics: Laws and Properties}. Oxford University Press.
\bibitem[Fine(1994)]{Fine1994} Fine, K. (1994). Essence and modality. \emph{Philosophical Perspectives}, 8, 1--16.
\bibitem[Kripke(1980)]{Kripke1980} Kripke, S.~A. (1980). \emph{Naming and Necessity}. Harvard University Press.
\end{thebibliography}

\end{document}
